\BeleskeID{09_3-s056}
% element: 09_3-s056:e05
% slide-id: 09_3-s056
Na izlaz prolazne logike obično se dodaje standardno CMOS kolo, odnosno invertor, da sniženi naponski nivoi ne bi nastavili da se prenose. Tako se mogu dobiti i NI ili NILI funkcije.

Dodatni tranzistor $R$ služi restauraciji visokog nivoa unutrašnjeg čvora. Kada prolazna mreža daje jedinicu, izlaz invertora je nula. Ta nula uključi P tranzistor $R$, koji dopuni parazitne kapacitivnosti do punog $V_{DD}$.

Pri prelasku izlaza invertora sa nule na jedinicu prolazna mreža mora povući njegov ulaz dovoljno nisko da N tranzistor prestane da provodi. Tranzistor $R$ se zato dimenzioniše tako da ne spreči to povlačenje. Situacija je analogna određivanju niskog izlaznog nivoa pseudo-NMOS invertora; važe isti odnosi strujnih sposobnosti.
